\subsection{The transfer eigenvalue on a finite circle}
\label{subsec:bethe-transfer}

We now set $c=2$ and $\mathcal W=W$.  Proposition~\ref{prop:bethe-ground}
identifies the solution of \eqref{eq:B8} on which the transfer norm is
attained.  We must still calculate that eigenvalue for the heat kernel on
the circle.  The rotor representation carried across the seam contributes
a sum which is absent from a calculation with infinite-line plane waves.
We will first obtain this sum exactly and then bound its nontrivial terms
uniformly when $n_p$ grows with $W$.

Write $n_p=2M+\ell$, with $\ell\in\{0,1\}$, and retain the twists
$\zeta,\chi$ used in \eqref{eq:B8}.  We conjugate each exterior binary
index by the antisymmetric tensor.  This fixed change of basis reverses
the diagonal matrix-index twists.  The seam therefore acts in rotor
sector $(q,j)$ by the scalar $e^{i\zeta q}$ and by $g_j^{-1}$ on each of
its two matrix indices, where $g_j$ is the spin-$j$ representation of
$g=\operatorname{diag}(e^{i\chi},e^{-i\chi})$.  Throughout this calculation
\[
 q\in\mathbb Z,\qquad j\in\tfrac12\mathbb Z_{\ge0},\qquad
 2j\equiv q\pmod 2,\qquad d_j=2j+1,\qquad
 E_{qj}=\tfrac12(q^2+d_j^2-1).
\]
The normalized Haar harmonics are $\sqrt{d_j}\,D^j_{ab}(Q)e^{iqv}$.
By Lemma~\ref{lem:brownian-ground}, the simple balanced ground space of
\eqref{eq:B2} is invariant under the transfer operator.  Near zero twist
its analytic continuation is still one-dimensional.  Denote the transfer
scalar on that space by $\Lambda_{n_p}(\zeta,\chi)$; at zero twist it is
$\|T_{n_p}\|$.

\begin{lemma}[Finite-circle transfer formula]
\label{lem:bethe-transfer-formula}
For the root branch of Lemma~\ref{lem:bethe-roots}, with fixed finite
$n_p$ and $W$, let $k_t$ be its momenta and let $t_j$ be the auxiliary
trace eigenvalue in \eqref{eq:B9}.  The function
$\Lambda_{n_p}(\zeta,\chi)$ is given near zero twist by
\begin{equation}
 \Lambda_{n_p}(\zeta,\chi)
 =\lambda^{n_p}\sum_{q,j}e^{-WE_{qj}+i\zeta q}
 \frac{(-2)^{n_p}i^{-2n_p}t_j(-iq)^2}
 {\displaystyle\prod_{t=1}^{n_p}
   \bigl[(d_j+1)^2-(ik_t-q)^2\bigr]
   \bigl[(d_j-1)^2-(ik_t-q)^2\bigr]} .
 \label{eq:B20}
\end{equation}
The sum is over the admissible rotor sectors just specified.  Apparent
singularities in this expression are interpreted by analytic
continuation from parameters at which coefficient extraction and the
Bethe vectors are nonsingular.
\end{lemma}

\begin{proof}
At an ordinary point of the root branch, choose a momentum ordering
whose color coefficient $A$ in the Bethe wave is nonzero.
The transfer preserves its eigenline, so the coefficient of
$\exp(i\sum_a k_aY_a)$ in the transferred wave is
$\Lambda_{n_p}A$.  We compute this coefficient, initially at the
nonresonant parameters specified below.

We first compute the contribution near one output coordinate $Y$.  Let
$X$ be the integration coordinate assigned to it, let $k$ be the momentum
carried by $X$, and fix the rotor sector $(q,j)$ on the intervals outside
these two insertions.  Put $d=2j+1$ and $z=ik-q$.  If $S_j$ denotes the
spin-$j$ angular momentum, the operators
$\eta_L=S_j\cdot\sigma$ and $\eta_R=S_j\cdot\sigma$ act on the left and
right rotor and binary indices, respectively.  These two operators
commute.  We extract the coefficient of $e^{ikY}$ obtained by taking
$Y$ as the integration endpoint for $X$.  The sum of this coefficient
over the two orders $X<Y$ and $X>Y$ is
\begin{equation}
 \frac{-2(z-2\eta_L)(z-2\eta_R)}
 {\bigl((d+1)^2-z^2\bigr)\bigl((d-1)^2-z^2\bigr)}.
 \label{eq:B19}
\end{equation}
We verify both the normalization and the endpoint denominators in this
formula.

Suppose first that $X<Y$, with no other insertion between them.
Multiplication by the conjugate entry at $X$ changes the intermediate
sector to $(q-1,j\pm\tfrac12)$.  Its energy minus $E_{qj}$ is
$1\pm d-q$.  On one matrix index the Clebsch--Gordan projections in
$j\otimes\tfrac12$ are
\[
 P_+(\eta)=\frac{j+1+\eta}{d},\qquad
 P_-(\eta)=\frac{j-\eta}{d}.
\]
The two matrix indices supply $P_\pm(\eta_L)P_\pm(\eta_R)$, with the same
choice of sign because they belong to the same intermediate rotor
representation.  The normalization of the Haar harmonics supplies the
additional factor $d/(d\pm1)$.  Indeed, multiplication from a normalized
spin-$j$ matrix element to a normalized spin-$(j\pm\tfrac12)$ matrix
element has coefficient $\sqrt{d/(d\pm1)}$ times the two
Clebsch--Gordan coefficients.  The return insertion is its adjoint and
contributes the same factor.  Thus this
dimension factor occurs once in the product of the two insertions, not
once for each matrix index.

For $X<Y$, the input plane times the heat factor is
$e^{ikX}e^{-(1\pm d-q)(Y-X)}$.
Its upper-endpoint contribution is $e^{ikY}/(1\pm d+z)$.
For $X>Y$ the lower-endpoint calculation gives denominator
$1\pm d-z$.  The exterior transpose and the conjugation by the
antisymmetric tensor replace $\eta_L,\eta_R$ by $-\eta_L,-\eta_R$ in
that contribution.  Consequently the full coefficient is
\[
 \sum_{s=\pm1}\frac{d}{d+s}
 \left\{
 \frac{P_s(\eta_L)P_s(\eta_R)}{1+sd+z}
 +\frac{P_s(-\eta_L)P_s(-\eta_R)}{1+sd-z}
 \right\}.
\]
Substituting the two projection formulas and taking a common denominator
gives \eqref{eq:B19}.  If $d=1$, the minus representation is absent:
we omit it before simplifying.  The result is the canceled value of
\eqref{eq:B19}, in particular $2/(4+k^2)$ when $(q,j)=(0,0)$.

We next extract the coefficient of a full output plane
$\exp(i\sum_a k_aY_a)$ in the actual heat-kernel integral.
Initially choose a nonsingular Bethe vector and parameters for which
all integrated partial-sum rates are nonzero, no $ik_a$ is a pure
rotor-energy difference, and no $i(k_a-k_t)$, for $a\ne t$, is such a
difference.  The relevant energy differences are even integers.
These conditions can be met on the branch: a generic common twist
$\zeta$ removes pure-rotor coincidences and zero partial-sum rates.
Near the large-circle end, the separated centers in \eqref{eq:B18}
and an infinitesimal spin twist $\chi$ resolve the remaining
pair-difference resonances.  Here $n_p$ and $c=2$ are fixed while the
circumference varies.  We will pass to all other parameters by analytic
continuation after computing the coefficient.

Fix a total ordering of $X$ and $Y$ and expand in rotor harmonics.
An ordered block $x_1<\cdots<x_r$ of input coordinates in a gap
$(a,b)$ between consecutive output coordinates has exponential rates
$\alpha_1,\ldots,\alpha_r$.  Repeated integration produces terms with
endpoint phases
\[
 \exp\left\{a\sum_{h=1}^{s}\alpha_h
              +b\sum_{h=s+1}^{r}\alpha_h\right\},
 \qquad 0\le s\le r.
\]
Thus a prefix is assigned to the left endpoint and the remaining
suffix to the right; this description includes the nested integrations
whose intermediate endpoints coincide.  Over all gaps it gives a
weakly order-preserving assignment of inputs to outputs or to the
seam endpoints $0,W$.

At an output $Y_a$, the exponent is the sum of the assigned input
momenta, multiplied by $i$, plus the adjacent rotor-energy difference.
If no input is assigned there, the exponent is purely a rotor-energy
difference and cannot equal $ik_a$.  Every output therefore receives
at least one input.  Since their numbers agree, exactly one input is
assigned to each output and none to the seam.  The difference
nonresonance then forces that input to carry $k_a$, and forces its
remaining rotor-energy difference to vanish.

In particular, a contributing gap contains at most two inputs, one
assigned to each endpoint.  When two integration variables lie in the
same $Y$ gap and are assigned
to its opposite ends, the coefficient is the product of their two
endpoint factors.  This can be checked without separating the ordered
integrations: for exponential rates $\alpha,\beta$, the term assigning
$x_1$ to $a$ and $x_2$ to $b$ in
\[
 \int_a^b e^{\beta x_2}\int_a^{x_2}e^{\alpha x_1}\,dx_1\,dx_2
\]
has coefficient $-1/(\alpha\beta)$, precisely the lower and upper
endpoint coefficients multiplied together.  The other terms assign the
two variables to a common endpoint and do not contribute to the required
plane.  After every $X$ has been paired with a $Y$, the vanishing
energy differences show that all intervening intervals carry the
same rotor energy.  The charge has already returned to the same $q$,
so equality of energies forces the same $j$.
Thus one fixed through sector $(q,j)$ contributes the ordered product of
\eqref{eq:B19} over the momenta, followed by its seam trace.

To justify coefficient extraction also for the infinite rotor sum,
work in the ordered complex tube
$0<\operatorname{Re}Y_1<\cdots<\operatorname{Re}Y_{n_p}<W$.
At fixed finite $n_p$, intermediate labels differ from the through
labels by bounded steps, so their energies differ from $E_{qj}$ by
$O_{n_p}(1+|q|+d_j)$.  On each compact subset of this tube, the
endpoint terms are therefore bounded by a fixed polynomial in the
labels times
\[
 \exp\{-WE_{qj}+C(1+|q|+d_j)\}.
\]
The quadratic energy gives normal convergence of the analytic
exponential sums, including after any finite imaginary translation.
Modulo even-integer shifts, their exponent vectors have only finitely
many types.  The commuting translations $Y_a\mapsto Y_a+i\pi$
multiply an exponential of rate $\gamma_a$ by $e^{i\pi\gamma_a}$;
these multipliers distinguish the types.  Polynomial interpolation in
the translations consequently projects onto each type.  In the type
containing the required plane, divide by $\exp(i\sum_a k_aY_a)$.
The resulting locally convergent Laurent series in $w_a=e^{2Y_a}$,
on $1<|w_1|<\cdots<|w_{n_p}|<e^{2W}$, has the desired plane
coefficient as its constant coefficient.  Uniqueness
of Laurent coefficients proves that no other term of the infinite
sum changes the extraction above.

Finally,
\[
 i(z-2\eta)=(-iq-k)-2iS_j\cdot\sigma
           =\mathcal L_j(-iq-k).
\]
The ordered product in each matrix index is therefore the monodromy
whose twisted trace has eigenvalue $t_j(-iq)$ in \eqref{eq:B9}.
The two indices give its square.  Multiplying by $\lambda^{n_p}$ and
summing the seam weights yields \eqref{eq:B20}.  At zero twist this is
the top transfer eigenvalue by Proposition~\ref{prop:bethe-ground} and
Lemma~\ref{lem:brownian-ground}; exceptional ring lengths follow by
continuity.
\end{proof}

\subsection{Uniform control of the nontrivial rotor sectors}
\label{subsec:bethe-winding}

We now estimate \eqref{eq:B20} in the regime of
Proposition~\ref{prop:brownian-parity}.  Its $(q,j)=(0,0)$ term is
\[
 Z_{n_p}(\zeta,\chi)
   :=\prod_{t=1}^{n_p}\frac{2\lambda}{4+k_t(\zeta,\chi)^2}.
\]
The issue is to compare every other term with this product without
accumulating a fixed loss for each of the $O(Wb^2)$ particles.

\begin{lemma}[Uniform transfer asymptotics]
\label{lem:bethe-transfer-uniform}
Fix $A,D>0$ and positive constants $c_-,c_+$.  Suppose $b\to\infty$,
\[
 c_-e^{\pi b}b^{-1/2}(\log b)^D
 \le W\le c_+e^{\pi b}b^{-1/2}(\log b)^D,
 \qquad |n_p/W-b^2|\le Ab.
\]
There is $c'>0$, depending only on these fixed constants, such that at
zero twist the sum in \eqref{eq:B20} over $(q,j)\ne(0,0)$, divided by
$Z_{n_p}(0,0)$, is $O(e^{-c'W})$, uniformly over these particle numbers.
If $u_r$ are the real centers of \eqref{eq:B13} at $\mathcal W=W$, then
\begin{equation}
 \log\|T_{n_p}\|
   =\sum_{r=1}^M p(u_r)+\ell f(0)+o(W^{-1}),
 \label{eq:B21}
\end{equation}
with the error uniform over the same band.
\end{lemma}

\begin{proof}
Set the spin twist to zero, initially retaining continuation from
arbitrarily small nonsingular spin twists where needed.  In
\eqref{eq:B9}, evaluated at $-iq$, every shifted argument of $Q_0$
has an odd integral imaginary part relative to a real formal root.
Indeed $2m\equiv2j\equiv q\pmod2$, so $q+2m\pm1$ is odd, and the
same is true of $q\pm(2j+1)$.  These factors are uniformly separated
from zero.

At the formal pair spacing $k_{r\pm}=u_r\pm i$, the contribution of
this pair to $F$ and its accompanying $Q_0$ factors in each summand of
\eqref{eq:B9} cancels to
\[
 (-iq-u_r)^2+d^2,\qquad d=2j+1,
\]
independently of $m$.  It follows that, in every one of the $d^2$
summands obtained by squaring $t_j$, the absolute contribution of this
pair relative to its $(0,0)$ contribution is
\[
 \frac{(u_r^2+1)(u_r^2+9)}
 {\left|\bigl((iu_r-q)^2-(d+2)^2\bigr)
              \bigl((iu_r-q)^2-(d-2)^2\bigr)\right|}
 \le1.
\]
To see the inequality, factor the denominator into four linear factors.
Their absolute values are $\sqrt{u_r^2+w^2}$, where the four widths are
$|q\pm(d+2)|$ and $|q\pm(d-2)|$.  All are positive odd integers, and
at least two are at least three.  Their product is consequently at
least $(u_r^2+1)(u_r^2+9)$.

This comparison remains uniform for the true roots.  The linear factors
just used, and the individual factors before cancellation for the two
pair constituents, stay a fixed distance from zero.  Equations
\eqref{eq:B16}--\eqref{eq:B17} make the root corrections smaller than
every power of $W^{-1}$.  There are only polynomially many factors in
$W$ in the particle band.  Hence the total multiplicative change is
bounded, and in fact differs from one by less than every prescribed
power of $W^{-1}$.  This is the point at which an estimate per particle
would have been insufficient without the sign-free bound of one at the
formal spacing.

If $\ell=1$, denominators contributed by the single momentum can vanish
at zero common twist before cancellation.  We control them on one
complex $\zeta$ circle instead.  Choose a sufficiently large fixed
$C$ and put $|\zeta|=W^{-C}$.  The common twist translates every
momentum and auxiliary root by $\zeta/W$.  At $\zeta=0$ the single
momentum is zero up to an error smaller than every power of $W^{-1}$.
Thus on this circle every possibly vanishing single factor has modulus
bounded below by a fixed multiple of $W^{-C-1}$.  All single-root
contributions are consequently bounded by a fixed power of
$1+W+d+|q|$.  The change in all pair factors together is bounded as
well: each argument moves by $W^{-C-1}$ and the number of factors is
polynomial in $W$.  Choosing $C$ once sufficiently large controls their
product.

We now justify using the maximum-modulus principle on this circle;
this requires analyticity on its whole interior, with no appeal to a
lower bound for the gap between different total-momentum sectors.
Let $\psi_0$ be the normalized untwisted ground vector, with total
momentum zero, and define the change of gauge
\[
 (U_\zeta\psi)(X)
 =\exp\!\left(\frac{i\zeta}{W}\sum_{t=1}^{n_p}X_t\right)\psi(X).
\]
After this change of gauge, the Hamiltonian on periodic functions is
$H_2+2\zeta P/W+n_p\zeta^2/W^2$, where
$P=-i\sum_t\partial_{X_t}$.  In the $P=0$ subspace its ground line is
therefore the same simple line for every $\zeta$.
Write $T_{n_p}(\zeta,0)$ for the twisted transfer.  Its rotor expansion
is analytic, and so is the fixed-gauge scalar
\[
 \left\langle\psi_0,
 U_\zeta^{-1}T_{n_p}(\zeta,0)U_\zeta\psi_0\right\rangle.
\]
Near zero this scalar equals $\Lambda_{n_p}(\zeta,0)$ by commutation.
It therefore continues $\Lambda_{n_p}$ across any crossings with
other momentum sectors.  The coefficient formula may first be used
with arbitrarily small spin twists and then continued to zero spin
twist, equivalently taking the removable analytic limits in
\eqref{eq:B20}.

On the $\zeta$ circle, the ratio of the nontrivial-sector sum to
$Z_{n_p}$ is therefore bounded by
\[
 C\sum_{(q,j)\ne(0,0)}
 d^2(1+W+d+|q|)^K
 \exp\{-WE_{qj}+W^{-C}|q|\}
 =O(e^{-c'W})
\]
for fixed $K$ and some $c'>0$.  Every nontrivial admissible sector has
$E_{qj}\ge2$, and the energies grow quadratically in $q,d$, so the
polynomial factors and the number of summands are absorbed by the
exponential.  The ratio is analytic on the enclosed disk:
$Z_{n_p}$ is analytic and nonzero there, and the whole nontrivial-sector
sum is $\Lambda_{n_p}-Z_{n_p}$.  The maximum-modulus principle gives
the same bound at zero twist.  For $\ell=0$ the separated pair factors
already give the bound directly; the same circle argument also applies.

It remains to take the logarithm of the leading product.  Since
$2\lambda=b^2$, a formal pair contributes
\[
 \frac{b^4}{(4+(u_r+i)^2)(4+(u_r-i)^2)}
 =\frac{b^4}{(u_r^2+1)(u_r^2+9)}=e^{p(u_r)},
\]
and the possible zero single contributes $b^2/4=e^{f(0)}$.
The cumulative error from replacing the true roots by the centers of
\eqref{eq:B13} is smaller than every power of $W^{-1}$ by
\eqref{eq:B17}; the nontrivial-sector error is exponentially smaller.
The positive zero-twist eigenvalue therefore has the logarithm
\eqref{eq:B21}.
\end{proof}

\subsection{Maximizing over the two particle parities}
\label{subsec:bethe-parity-max}

The transfer asymptotics have reduced the problem to a sum over real
pair centers.  The dressed energy $\epsilon$ of \eqref{eq:B3}, with
endpoint $B=B_*$, expresses this sum as a common bulk term, an explicit
single-particle cost, and a quadrature error.  Recall that
$\epsilon_{\rm in}=\epsilon\mathbf1_{(-B_*,B_*)}$ is also the positive
part of $\epsilon$.

Let $d(x)$ now denote the globally defined counting derivative for
\eqref{eq:B13} at $\mathcal W=W$; it is not truncated to its counting
interval $(-B_c,B_c)$.  Summing $p=\epsilon-K*\epsilon_{\rm in}$ at
the centers and using the definition of $d$ gives the exact identity
\begin{equation}
 \sum_{r=1}^M p(u_r)+\ell f(0)
 =\frac{W}{\pi}\int_{\mathbb R}\epsilon_{\rm in}(x)\,dx-\ell m_*
  +\left[\sum_{r=1}^M\epsilon(u_r)
        -W\int_{\mathbb R}\epsilon_{\rm in}(x)d(x)\,dx\right].
 \label{eq:B22}
\end{equation}
Indeed,
\[
 W\int\epsilon_{\rm in}d
 =\frac{W}{\pi}\int\epsilon_{\rm in}
   +\sum_{r=1}^M(K*\epsilon_{\rm in})(u_r)
   -\ell\bigl((K_1-K_3)*\epsilon_{\rm in}\bigr)(0),
\]
and the remaining single term is $-\ell m_*$ by its definition.

\begin{lemma}[The parity variational error]
\label{lem:bethe-parity-error}
Under the parameter assumptions of Lemma~\ref{lem:bethe-transfer-uniform},
the bracket in \eqref{eq:B22} is at most $C/W$, uniformly in the particle
band.  In each parity there is a particle number in a band of the same
form for which the bracket is at least $-C/W$.  Specifically, one may
take $M$ to be a nearest integer to $W\int\rho_{B_*}$ and set
$n_p=2M+\ell$.
\end{lemma}

\begin{proof}
For the upper bound, replace each $\epsilon(u_r)$ by its positive part.
Use the counting coordinate
\[
 s(x)=\int_0^x d(y)\,dy,
 \qquad s(u_r)=\frac{r-(M+1)/2}{W},
 \qquad s(\pm B_c)=\pm\frac{M}{2W}.
\]
The centers are exactly the midpoints of cells of length $1/W$ in this
coordinate.  The estimates accompanying \eqref{eq:B14} give $d\ge c>0$
and polynomial bounds in $b$ for its fixed-order derivatives on the
needed interval.  The integral equation \eqref{eq:B3} gives such bounds
for the fixed-order derivatives of $\epsilon$.  Moreover
$|\epsilon'|\le\pi/2$ on the whole line by
Proposition~\ref{prop:brownian-asymptotics}.  Thus, for
the inverse function $x=x(s)$,
\[
 \left|\frac{d}{ds}\epsilon(x(s))\right|
 =\left|\frac{\epsilon'(x(s))}{d(x(s))}\right|\le C,
\]
while each further fixed-order derivative is polynomially bounded in
$b$.

Set $g(s)=\epsilon_+(x(s))$.  On any consecutive collection of full
cells in a smooth portion, with outer endpoints $\alpha,\beta$, the
composite midpoint expansion gives
\[
 \sum_{\text{these cells}}g(s_r)
       -W\int_\alpha^\beta g(s)\,ds
 =-\frac{g'(\beta)-g'(\alpha)}{24W}
       +O\!\left(\frac{\operatorname{poly}(b)}{W^2}\right).
\]
The first derivatives are uniformly bounded, and the higher derivatives
and interval length cost only polynomial factors in the remainder.
Thus this error is $O(W^{-1})$ with a constant independent of $b$.
At most two cells meet the cutoffs
$x=\pm B_*$.  Since $\epsilon$ vanishes there and its derivative in
$s$ is bounded, these cells contribute $O(1/W)$ as well.  Consequently
\[
 \sum_r\epsilon(u_r)
 \le W\int_{-B_c}^{B_c}\epsilon_+(x)d(x)\,dx+\frac{C}{W}
 \le W\int_{\mathbb R}\epsilon_{\rm in}(x)d(x)\,dx+\frac{C}{W}.
\]
The last inequality uses positivity: any positive portion outside the
counting interval is only omitted on its left-hand side.

For the lower bound choose $M$ as in the statement.  Equation
\eqref{eq:B5} and the relation between $b$ and $B_*$ imply
$2\int\rho_{B_*}=b^2+O(b)$, so these particle numbers satisfy the
band condition after enlarging its fixed constant.  Since $\int_{-B_c}^{B_c}d=M/W$,
\eqref{eq:B14}, including its uniform $L^1$ bound for the impurity
correction, yields
\[
 \left|\int\rho_{B_c}-\int\rho_{B_*}\right|=O(W^{-1}).
\]
The densities increase with their interval: this follows by positive
iteration of their defining equation.  Their difference in total mass
therefore bounds from below the mass on the excess or missing interval.
First $\rho_B\ge1/\pi$ gives $|B_c-B_*|=O(W^{-1})$.  Those intervals
then lie within distance one of the relevant larger endpoint; applying
the edge bound $\rho_B\ge c\sqrt B$ from \eqref{eq:B5} improves this to
\[
 |B_c-B_*|\le\frac{C}{W\sqrt{B_*}}.
\]
On this excess or missing interval the counting density obeys
$d\le CB_*$, by \eqref{eq:B14} and its defining convolution, and
$|\epsilon(x)|\le C\,\bigl||x|-B_*\bigr|$.  Therefore replacing the
counting interval by $(-B_*,B_*)$ in the dressed-energy integral costs
at most
\[
 CWB_*|B_c-B_*|^2\le\frac{C}{W}.
\]
The same midpoint estimate, now for the smooth function $\epsilon$
on the counting interval, compares its sum and integral with error
$O(1/W)$.  Combining these two bounds proves the lower estimate for the
bracket.
\end{proof}

We can now complete Proposition~\ref{prop:brownian-parity}.  By its
particle-band hypothesis, every maximizing particle number is covered
by the uniform upper bound above.  After one fixed enlargement of the band constant, the choices in
Lemma~\ref{lem:bethe-parity-error} are covered as well and give the
matching lower bounds.  Thus
\[
 \log\tau_\ell
 =\frac{W}{\pi}\int\epsilon_{\rm in}-\ell m_*+O(W^{-1}),
 \qquad \ell=0,1.
\]
By \eqref{eq:B5},
\[
 e^{-\pi B_*/2}=4\pi e B_*e^{-\pi b}(1+o(1)),
 \qquad
 m_*\sim32\sqrt{B_*}\,e^{-\pi b}
       \sim32\sqrt2\,b^{1/2}e^{-\pi b}.
\]
The prescribed size of $W$ gives $Wm_*\asymp(\log b)^D\to\infty$.
Subtracting the two parity formulas gives
\[
 \log\tau_0-\log\tau_1=m_*+O(W^{-1}),
\]
and therefore proves \eqref{eq:B1}.
